\section{Data：数据创造与能力构建}

\subsection{新 Domain = 新 Capability}
讲师反复强调 K2.5 并非在旧模型上贴标签，而是认为 ``新的 domain 就意味着新的能力切片''：图文、API、GUI、桌面工具乃至 OpenClaw 控制台，都被设计成各自的 benchmark。每一个 domain 的数据采集、标注规范和 success criteria 都决定了 agent 最终要学会哪类行为。我们阅读技术报告时的首要问题就是：这个 domain 的数据是由谁定义、作何合成、给出了什么样的输出？

早期的混合策略只在 10\% 以内加入图像（1:9），到了中期 250k steps 变成 2:8，最后变成 1:1。但讲师关键一句话是：``只要在 Step0 引入视觉，文本与编程能力就会同步走高''，说明视觉并不是附加，而是塑造语言理解的核心维度。

\begin{importantbox}{早期混合视觉的杠杆}
在训练初期只投放 10\% 左右的图文数据（1:9）、middle 走到 20\%（2:8），到后期才到 1:1。讲师的数据表明只要从 Step0 就混入视觉就能显著提升 \textbf{文本} 和 \textbf{编程} 能力，说明 K2.5 的视觉不是附加项目，而是在语言表示里不断灌入新 signal。
\end{importantbox}

\subsection{合成数据方法}
K2.5 的数据系统高度可复现：把目标能力拆成 benchmark（如屏幕指令、跨界回答、长文推理），为每个 benchmark 设计数据生成 pipeline，再用 SFT + RL 的组合训练 baseline 检查数据质量。讲师强调我们要把每个 benchmark 的 ``完成标准'' 写清楚，这样才能看出模型什么时候做完 task。

常规流程依旧是 1) 定义 benchmark，2) 用规则/模版生成样本，3) 让 baseline 验证、再进入 RL。这个 pipeline 在 lecture 22:00~32:00 的章节中被反复提到，特别是在 video-to-code / GUI 合成样本时会把 API trace 与 UI 状态同时写入 prompt。

\begin{enumerate}
    \item \textbf{定义 benchmark}：指定几类任务（如 video-to-code、GUI 操作、长文理解），厘清 success criteria 和可视化指标。
    \item \textbf{设计算法生成 pipeline}：用图像检索、tool simulator、prompt-driven grammar 生成合成样本，并记录每一步的 metadata。
    \item \textbf{baseline 验证}：用 SFT 模型跑 benchmark，确认 reward 侧的仿真分数才允许进入 RL。
\end{enumerate}

\begin{table}[H]
\footnotesize
\centering
\begin{tabular}{p{3cm}p{4cm}p{7cm}}
\toprule
训练阶段 & 视觉 / 文本 比例 & 观察到的效果 \\
\midrule
初期（Step0） & 1:9 & 虽然视觉样本少，但只要从第一步就混入，文本和代码表现就迅速抬头，短期内便能在多个 benchmark 上产生 \textbf{知识增强}。 \\
中期（Step 250k） & 2:8 & 视觉占比提升，让中间层学会在 sequence 中插入工具调用、GUI 操作与场景感知，缓解了多模态梯度冲突。 \\
后期（Step 末端） & 1:1 & 视觉与文本并重，使最终评估的 critical steps 同时考察视觉推理与语言执行，提升 agentic performance。 \\
\bottomrule
\end{tabular}
\caption{K2.5 在不同阶段控制的图文混合比例与训练效果，数据来自讲师附录 Figure 9 的曲线。}
\end{table}

\subsection{Domain-specific synthesis loops}
每个 benchmark 都有自己的构建细节：video-to-code 要把 UI/landing page、API 端点、设计稿切成多帧并对齐文本；长文理解需要做 layout-aware OCR；GUI/CLI 的 sample 里要保留 tool trace、模拟点击、记录 API 响应。讲师强调，这些数据 pipeline 直接决定了 agent 的 prompt 结构与生成长度，因而成为 ``domain-to-capability'' 的关键环节。我们在 lecture 15:40~18:10 看到，讲师用``Tool trace + UI 状态''两列表格，把每个 sample 里要记录的字段列出。

\begin{knowledgebox}{Domain-specific synthesis 的三条策略}
\begin{itemize}
    \item 把每个 domain 的 benchmark 拆成 intent、tools、output 三段，明确何时需要 GUI、何时需要 API、何时只需要 text。
    \item 用 tool simulator 、UI mock generator 生成带 trace 的样本，保证训练期间有一致的 token budget。
    \item 用 uniform metadata schema 记录 timestamp、frame id、tool log，便于 later evaluation 回放每一个 critical step。
\end{itemize}
\end{knowledgebox}

\subsection{质量治理与可解释}
讲师指出：数据 pipeline 不能只关注数量，还要把 quality 控制在 ``critical steps'' 线内。每批样本在 SFT baseline 上跑一轮，如果 reward 分数低于阈值，就会回到数据模块重新校准 prompt grammar。此外，visual sample 会附上 frame-level label，方便 later debugging 时追溯每一个失误的视觉 Token。

\begin{importantbox}{数据质量的自动化 Guardrail}
K2.5 用两个层级的质量检查：第一层是 SFT baseline 的 reward gate，第二层是 human-in-the-loop 检查 GUI 操作和屏幕文字。只有同时通过两层的 sample 才能进入 RL，避免了 ``reward hacking'' 的爆炸式噪音。
\end{importantbox}

\subsection{本章小结}
K2.5 的 data story 是 ``domain-to-capability'' 的映射：用 video-to-code、对话推理、GUI/工具三个 benchmark 类别串联数据能力，并通过 SFT baseline + critical steps 维持质量，使得训练样本始终保持可解释的 metadata。
